
\documentclass[11pt]{article}

% --- Packages
\usepackage[margin=1in]{geometry}
\usepackage{amsmath,amssymb,mathtools,bm}
\usepackage{amsthm}
\usepackage{graphicx}
\usepackage{subcaption}
\usepackage{booktabs}
\usepackage{siunitx}
\usepackage{microtype}
\usepackage{enumitem}
\usepackage{framed}
\usepackage{hyperref}
\hypersetup{colorlinks=true, allcolors=blue}

% --- Title
\title{PER as an Information--Dynamical Principle:\\
Potential--Evanescent--Real, Variational Structure, and Data Tests\\
{\large with LIGO GW150914 and the NIST CH Pockels Bell Test}}
\author{}
\date{\today}

% --- Macros
\newcommand{\R}{\mathbb{R}}
\newcommand{\E}{\mathbb{E}}
\newcommand{\p}{\bm{p}}
\newcommand{\f}{\bm{f}}
\newcommand{\one}{\bm{1}}
\newcommand{\inner}[2]{\left\langle #1,#2 \right\rangle}
\newcommand{\brak}[1]{\left( #1 \right)}
\newcommand{\Brak}[1]{\left[ #1 \right]}

% Fisher/Shahshahani metric
\newcommand{\fish}[2]{\sum_i \frac{#1_i\,#2_i}{p_i}}
\newcommand{\normp}[1]{\left\| #1 \right\|_p}

% Theorems
\newtheorem{lemma}{Lemma}
\newtheorem{prop}{Proposition}
\newtheorem{remark}{Remark}

\begin{document}
\maketitle

\begin{abstract}
We present \emph{PER}---a Potential$\to$Evanescent$\to$Real pipeline---that unifies selection, intrinsic time, and recording into one compact formalism.
The evanescent stage is a Fisher--metric (Shahshahani) replicator ascent of a coherence potential $C(\p,t)$; the record is a stopping time when a concentration or coherence threshold is crossed.
We derive Onsager--Machlup, Hamiltonian, and contact/GENERIC formulations; a Hamilton--Jacobi and time-reparametrized action; and diagnostics that forecast transient ``evanescence''.
We couple PER to an amplitude‑level SQC lift, obtaining a single Fisher-Bures variational principle that recovers weak‑measurement behavior; we illustrate with double‑slit interference and the Quantum Cheshire Cat, and we show how witness‑rate leads/nonunitarity‑leads lags align with rise/fall transitions in data.
We then test PER on two public datasets: (\textbf{I}) \textbf{LIGO GW150914}, where the witness ratio $R_w$ \emph{leads} the SQC nonunitarity derivative during the \emph{rise} by \SIrange{32}{48}{ms}, while a full-trace (rise+fall) statistic shows \emph{collapse-lead}; (\textbf{II}) the \textbf{NIST CH Pockels (100\,kHz) Bell test}, where conservative full-trace statistics show \emph{collapse-lead} with lags of \SIrange{7.6}{30}{s} depending on the SQC grid; early ROIs and more local witnesses shrink the magnitude.
Across both, PER's intrinsic clock and stopping rules are consistent, and $R_w,R_1$ and an FDR diagnostic spike at transitions.
\end{abstract}

\paragraph{Roadmap.}
Section~\ref{sec:per-core} derives the PER flow, core identities, and the intrinsic clock with stopping.
Section~\ref{sec:variational} develops Onsager--Machlup, Hamiltonian/contact, HJ, and time-reparametrized actions.
Section~\ref{sec:info-geom} gives the information geometry (Shahshahani), geodesics, integrability, and optimal-transport links.
Section~\ref{sec:per-sqc-unified} unifies PER and SQC; 
Section~\ref{sec:weak} views PER as weak-measurement; 
Section~\ref{sec:double-slit} illustrates double-slit interference; Section~\ref{sec:qcc} illustrates the Quantum Cheshire Cat.
% Section~\ref{sec:sqc} consolidates the Stochastic--Quantum Correspondence (SQC).
Section~\ref{sec:diagnostics} defines $R_1$, $R_w$, and an FDR diagnostic; Section~\ref{sec:design} records the experimental pipeline.
Sections~\ref{sec:gwosc}--\ref{sec:bell} present results for LIGO and Bell.
% Appendix~\ref{app:HK} treats Hellinger--Kantorovich; Appendix~\ref{app:kuramoto} summarizes Kuramoto scaling.

\paragraph{Notation}
\begin{center}
\begin{tabular}{@{}ll@{}}
\toprule
Symbol & Meaning (default) \\ \midrule
$p\in\Delta$ & pathway probabilities, $\sum_i p_i=1$ \\
$f_i=s_i+(Kp)_i-\lambda\,\kappa_i$ & PER drive (witness $s$, compatibility $K$, penalty $\kappa$) \\
$C(p,t)=\sum_i p_i f_i$ & coherence potential (Fisher gradient ascent) \\
$H(p)=-\sum_i p_i\ln p_i$ & Shannon entropy (PER) \\
$\theta$ & intrinsic clock, $\dot\theta=\alpha\,\dot C$ (forward/smoothed) \\
$\eta$ & MWU step (default $0.6$) \\
$\lambda$ & penalty weight (default $0$) \\
$R_1=-\dot H/\dot C$ & information–to–coherence ratio \\
$C_{\rm ext}=\sum_i p_i s_i$ & external witness contribution \\
$C_{\rm int}=p^\top K p-\lambda\sum_i p_i\kappa_i$ & internal contribution \\
$R_w=\dot C_{\rm ext}/(\dot C_{\rm int}+\varepsilon)$ & witness ratio (stabilized) \\
$\rho$ & SQC density operator with $\mathrm{diag}\,\rho=p$ \\
$Q_F=\|\Theta^\dagger\Theta-\mathbb I\|_F$ & SQC nonunitarity (dimensionless) \\
$d_{\rm uni}(\Gamma)$ & distance of step kernel to unistochastic set \\
$\alpha_\theta$ & clock scale in $\dot\theta = \alpha_\theta\,\dot C$ (units: time/evidence) \\
$\alpha_{\text{deph}}$ & coefficient in $\gamma(t)=\gamma_0+\alpha_{\text{deph}}\,[R_w(t)]_+$ (Sec.~\ref{sec:per-sqc}) \\
\bottomrule
\end{tabular}
\end{center}

% =========================================================
\section{PER: equations, identities, and intrinsic time}\label{sec:per-core}

\paragraph{Potential $\to$ Evanescent.}
Let candidates $i=1,\dots,N$ carry weights $p_i(t)\ge 0$ with $\sum_i p_i=1$.
Witness channels $(s(t),K(t))$ define
\begin{equation}
  f_i(t) \;=\; s_i(t) + \big(K(t)\,p(t)\big)_i - \lambda\,\kappa_i(t),
\end{equation}
and the \emph{evanescent} evolution is the Shahshahani replicator
\begin{equation}\label{eq:replicator}
  \dot p_i \;=\; p_i\big(f_i-\bar f\big),\qquad \bar f=\sum_j p_j f_j,
\end{equation}
the gradient ascent of $C(\p,t)=\sum_i p_i f_i(t)$ with respect to the Fisher metric
$\inner{u}{v}_p=\sum_i u_i v_i/p_i$:
\begin{equation}
  \dot \p \;=\; \mathrm{grad}_p\,C(\p,t).
\end{equation}

\noindent\emph{Assumptions.}
The identities below hold exactly when $s(t)$ and $K(t)$ are time–independent (or vary slowly so that coarse–graining with window $\Delta$ makes remainders negligible). When witnesses vary in time, an additive term $\dot C_{\rm ext}$ appears; we enforce $\dot\theta\ge 0$ via smoothing (Methods).

\paragraph{Boxed identities (quasi-static witnesses).}
When $(s,K)$ are frozen in time,
\begin{framed}
\begin{align}
  \dot C \;=\; \sum_i f_i\,\dot p_i \;=\; \sum_i p_i\,(f_i-\bar f)\,f_i \;=\; \mathrm{Var}_p(f) \;\ge 0, \label{eq:Cdot}\\[2pt]
  \dot H \;=\; -\sum_i \dot p_i\ln p_i \;=\;
  \sum_i p_i\,(f_i-\bar f)\,(-\ln p_i) \;=\; -\,\mathrm{Cov}_p\!\big(f,\ln p\big), \label{eq:Hdot}
\end{align}
\end{framed}
with $H=-\sum_i p_i\ln p_i$.
The information-to-coherence ratio
\begin{equation}
  R_1(t) \;:=\; \frac{-\dot H}{\dot C}
  \;=\; -\,\frac{\mathrm{Cov}_p(f,\ln p)}{\mathrm{Var}_p(f)}
  \;=\; -\,\mathrm{Corr}_p(f,\ln p)\,\frac{\sigma_{\ln p}}{\sigma_f},
\end{equation}
obeys $|R_1|\le \sigma_{\ln p}/\sigma_f$; it diverges at critical slowing-down when $\sigma_f\!\to\!0$.

\paragraph{Kernel term (correction).}
For time-independent witnesses,
\begin{equation}
 \frac{d}{dt}\Big(\tfrac12\,\p^\top K\,\p\Big)
 \;=\; \mathrm{Cov}_p\!\Big((Kp)_i,\ f_i\Big)
 \;=\; \mathrm{Var}_p\!\big((Kp)_i\big) \;+\; \mathrm{Cov}_p\!\Big((Kp)_i,\ s_i-\lambda\kappa_i\Big),
\end{equation}
reducing to a pure variance when the last term vanishes (e.g.\ $s_i-\lambda\kappa_i$ constant across $i$ or uncorrelated with $(Kp)_i$ under $p$).
When differentiating $\tfrac12 p^\top K p$, the identity $d/dt(\tfrac12 p^\top K p)=(Kp)\!\cdot\dot p$ assumes $K$ is symmetric.
For general $K$ use the symmetric part $(K+K^\top)/2$ in that derivative.


\paragraph{Discrete-time MWU identity.}
For multiplicative weights,
\begin{equation}
  p_i^{(t+1)} \;\propto\; p_i^{(t)}\,e^{\eta f_i^{(t)}},\qquad
  \ln p_i^{(t+1)} - \ln p_i^{(t)} \;=\; \eta f_i^{(t)} - \ln Z^{(t)},
\end{equation}
or, centered for intuition,
\begin{equation}
  \ln p_i^{(t+1)} - \ln p_i^{(t)} \;=\; \eta\big(f_i^{(t)}-\bar f^{(t)}\big)\;-\;\big[\ln Z^{(t)}-\eta\bar f^{(t)}\big],
\end{equation}
where the bracket is a normalization \emph{scalar} (nonzero except to $O(\eta)$).

\paragraph{Small–$\eta$ expansion.}
With $Z=\sum_k p_k e^{\eta f_k}=\mathbb{E}_p[e^{\eta f}]$, 
$\ln Z=\eta\,\bar f+\tfrac12\eta^2\mathrm{Var}_p(f)+O(\eta^3)$,
so the normalization bracket equals $\tfrac12\eta^2\mathrm{Var}_p(f)+O(\eta^3)$ (nonzero beyond first order).


\paragraph{Time as change: intrinsic clock and monotonicity.}
Define the PER clock $\theta$ by
\begin{equation}
  \dot\theta = \alpha\,\dot C \qquad \text{or} \qquad
  \theta(t)=\alpha\!\int_0^t\!\max\{\dot C(\tau),0\}\,d\tau,
\end{equation}
or by smoothing $C$ at scale $\tau$ so that $\dot\theta\ge 0$ at the chosen coarse-graining.
A \emph{record} forms when $p_k$ crosses a threshold (with Plackett--Luce tie-breaking).

\paragraph{When do we keep $\dot C\ge 0$?}
Sufficient conditions at a coarse scale: bounded total variation of $(s(t),K(t))$; mild time-scale separation (witnesses slower than selection); or explicit smoothing/averaging in $C$.

\paragraph{Units.} If $C$ is in “evidence–rate” units, choose $\alpha$ with units time/evidence so $\theta$ is dimensionless. 
For the forward–only clock we use a Gaussian or triangular kernel $\mathcal{S}_\tau$ (e.g.\ $\tau=\SI{50}{ms}$ for LIGO) to enforce $\dot\theta\ge0$ at the chosen coarse scale.

% =========================================================
\section{Variational, symplectic, contact, and HJ structure}\label{sec:variational}

\subsection*{Mapping between $C$ and the Hamiltonian; contact lift (dissipation explicit)}
The PER flow is a \emph{dissipative} Fisher–gradient ascent of $C$.
To place Hamilton–Jacobi notation alongside without erasing dissipation, set
\[
\mathcal{E}(p,t) := -\,C(p,t)+\text{const}(t),
\]
so the characteristic momentum $\pi_i=\partial_{p_i}W$ matches $\pi_i=f_i-\bar f$ on the \emph{slow manifold}.
Using a contact Hamiltonian $h=H-\gamma\,\mathcal{S}$ with $H(p,\pi,t)=\tfrac12\sum_i p_i\pi_i^2+\mathcal{E}(p,t)$ and 1–form $\eta=d\mathcal{S}-\sum_i\pi_i\,dp_i$, the dynamics contracts transversely (rate $\gamma>0$) onto
\[
\pi_i=f_i-\bar f,\qquad \sum_i p_i\pi_i=0,
\]
and the reduced flow on $\Delta^\circ$ is exactly $\dot p=\mathrm{grad}_pC$.
Thus the replicator is the \emph{dissipative} (contact) reduction of the lifted dynamics.

\subsection*{Onsager--Machlup / Rayleighian action}
On $T\Delta^\circ$ with Fisher metric $G(\p)$ and small noise $\varepsilon$,
\begin{equation}
\mathcal{A}_\varepsilon[\p(\cdot)] \;=\; \frac{1}{2\varepsilon}\int_0^T
 \big\|\dot\p - \mathrm{grad}_p\,C(\p,t)\big\|_p^2\,dt,
\end{equation}
whose extrema satisfy Eq.~\eqref{eq:replicator}.
Near criticality $\|\nabla C\|\!\to\!0$ the barrier collapses and large fluctuations proliferate.

\subsection*{Symplectic structure on $T^*\Delta$}
The cotangent bundle $T^*\Delta$ carries the canonical symplectic form $\omega = \sum_i dp_i\wedge d\pi_i$; in log-simplex fiber coordinates $\pi_i=\dot p_i/p_i$ with gauge constraint $\sum_i p_i\pi_i=0$ the Hamiltonian
\begin{equation}
H(\p,\bm{\pi},t) \;=\; \frac{1}{2}\sum_i p_i \pi_i^2 + \mathcal{E}(\p,t),
\qquad \pi_i=f_i-\bar f,
\end{equation}
generates the replicator characteristics.

\subsection*{Hamilton--Jacobi and Rayleighian link}
The principal function $W(\p,t)$ satisfies
\begin{equation}
  \partial_t W + \frac{1}{2}\sum_i p_i\left(\partial_{p_i}W\right)^2 + \mathcal{E}(\p,t) \;=\; 0.
\end{equation}
Its characteristics coincide with replicator trajectories; $W$ acts as a Lyapunov function whose level-sets track approach to recording.

\subsection*{Contact lift and GENERIC}
On the contact manifold with 1-form $\eta= d\mathcal{S} - \sum_i \pi_i\,dp_i$ and contact Hamiltonian $h=H-\gamma\mathcal{S}$, the slow manifold $\pi_i=f_i-\bar f$ is \emph{exponentially attracting} for $\gamma>0$; the gauge $\sum_i p_i\pi_i=0$ is enforced throughout.

\subsection*{Critical phenomena in the action}
For a generic one-parameter transverse crossing (saddle-node normal form), the action barrier scales as $\Delta\mathcal{A}\sim \|\nabla \mathcal{E}\|^2/\varepsilon$, predicting large fluctuations near criticality and the divergence of FDR defined below.

\subsection*{Time reparametrization invariance}
The intrinsic arclength $d\theta = \normp{\dot\p}\,dt$ yields a Jacobi--Maupertuis form
\begin{equation}
\mathcal{A}[\theta] \;=\; \int \left(\frac{1}{2}\left\|\frac{d\p}{d\theta}\right\|_p^2 - \mathcal{E}(\p,\theta)\right) d\theta,
\end{equation}
% and connects naturally to transport metrics in Appendix~\ref{app:HK}.

% =========================================================
\section{Information geometry, geodesics, and integrability}\label{sec:info-geom}

\paragraph{Shahshahani geometry (square–root chart).}
Equip $\Delta^\circ$ with the Fisher/Shahshahani metric $\langle u,v\rangle_p=\sum_i u_i v_i/p_i$ on $\sum_i u_i=0$.
Use the isometric square–root embedding
\[
x_i = 2\sqrt{p_i}\quad\Rightarrow\quad x\in \mathbb{S}^{N-1}_2\cap\mathbb{R}^N_{>0},
\]
so geodesics lift to great–circle arcs. This avoids coordinate artifacts and makes the gradient interpretation immediate; a concise derivation is given in Appendix~\ref{app:shahshahani}.

\paragraph{Geodesics and forcing.}
For quadratic $\mathcal{E}$ the replicator flow follows geodesics; deviations capture forcing by witnesses (``external fields'').

\paragraph{Integrability.}
For $N=2$ the flow is integrable; for $N\ge 3$ generic time-dependent forcing destroys integrability. Exact solutions persist for special $\mathcal{E}$ (e.g.\ quadratic/geodesic).

\paragraph{Optimal transport links.}
With the intrinsic time element, PER's evolution is equivalent to transporting probability along Fisher geodesics perturbed by a potential; Appendix~\ref{app:HK} details the mapping to unbalanced OT (HK).


% =========================================================
\section{PER/SQC combination and a unified action}
\label{sec:per-sqc-unified}

\paragraph{SQC layer (amplitude lift).}
Given PER marginals $p(t)\!\in\!\Delta$ and a classical step kernel $\Gamma(t\!\leftarrow\!0)$ with $p(t)=\Gamma p(0)$, the \emph{stochastic–quantum correspondence} (SQC) lifts $\Gamma$ to amplitudes $\Theta$ with
\[
|\Theta_{ij}(t\!\leftarrow\!0)|^2=\Gamma_{ij}(t\!\leftarrow\!0),\qquad
\sum_i|\Theta_{ij}|^2=1,
\]
and defines $\rho(t)=\Theta(t)\,\mathrm{diag}(p(0))\,\Theta^\dagger(t)$ so that $\mathrm{diag}\,\rho(t)=p(t)$.
Column phases are fixed by the \emph{minimal‑twist} gauge
$\min_\Theta\!\int\|\operatorname{skew}(\dot\Theta\,\Theta^\dagger)\|_F^2\,dt$ subject to $|\Theta|^2=\Gamma$.

\paragraph{Diagnostics on the SQC layer.}
We use $Q_F(t)=\|\Theta^\dagger\Theta-\mathbb{I}\|_F$ to quantify nonunitarity (disturbance/back‑action/coarse‑graining) and
$d_{\rm uni}(\Gamma)=\min_{U\in U(N)}\|\Gamma-|U|^2\|_F$ as a distance to unistochastic maps.

% --- In “Diagnostics on the SQC layer” ---
\paragraph{Scale and invariances.}
$Q_F=\|\Theta^\dagger\Theta-I\|_F$ is dimensionless and equals $0$ for column-orthonormal $\Theta$.
It is invariant to right-multiplication by diagonal phase matrices, and—when $|\Theta|^2=\Gamma$ exactly—to column permutations,
but increases under column rescalings that depart from unistochasticity. We therefore report $d_{\mathrm{uni}}(\Gamma)$ alongside $Q_F$.

\paragraph{Unistochasticity and approximation.}
A stochastic step $\Gamma$ admits $|\Theta|^2=\Gamma$ iff it is \emph{unistochastic} ($N=2$ always; not guaranteed for $N\ge3$).
When $d_{\rm uni}(\Gamma)>0$ we pick $\Theta$ in a minimal–twist gauge minimizing $\|\Gamma-|\,\Theta|^2\|_F$ subject to column normalization.
In that case $Q_F$ captures both physical nonunitarity \emph{and} factorization mismatch; we therefore report $d_{\rm uni}$ with $Q_F$ and treat it as a covariate in lag analyses.

\paragraph{Coupling PER and SQC (three standard modes).}
\begin{enumerate}[leftmargin=1.05em,itemsep=2pt,topsep=2pt]
\item \textbf{Passive lift (PER$\rightarrow$SQC).} Evolve $p$ by PER (replicator/MWU); factor the step kernel $\Gamma$; pick $\Theta$ with $|\Theta|^2=\Gamma$; set $\rho=\Theta\,\mathrm{diag}(p(0))\,\Theta^\dagger$. This provides phases/interference consistent with PER’s diagonals.
\item \textbf{Decoherence feedback (two‑way).} Let the PER witness ratio drive dephasing,
$\gamma(t)=\gamma_0+\alpha\,[R_w(t)]_+$,
and evolve
$\dot\rho=-i[H_{\rm eff}(t),\rho]+\sum_i\gamma_i(t)\big(P_i\rho P_i-\tfrac12\{P_i,\rho\}\big)$.
Optionally feed residual coherence back to PER by $f_i \leftarrow f_i + \beta \sum_{j\ne i}\mathrm{Re}[\rho_{ij}G_{ij}]$.
\item \textbf{Kernel‑consistent lift (tight coupling).} Discretize PER as $p\mapsto \Gamma p$ (e.g., softmax/Plackett–Luce or OT step) and \emph{require} $|\Theta|^2=\Gamma$ each step. This keeps $\mathrm{diag}\,\rho=p$ by construction.
\end{enumerate}

\paragraph{Unified variational principle.}
A single hybrid action on $(p,\rho,\Theta)$ recovers both flows as zero‑action limits and enforces consistency:
\[
\boxed{\;
\mathcal{S}[p,\rho,\Theta]=\!\int_0^T
\underbrace{\tfrac{1}{2\varepsilon}\big\|\dot p-\mathrm{grad}_p C(p,t)\big\|_{p}^{2}}_{\text{PER/Fisher}}
+\underbrace{\tfrac{1}{2\varepsilon_q}\big\|\dot\rho-\mathcal{L}_{p,t}(\rho)\big\|_{B}^{2}}_{\text{SQC/Bures}}
+\underbrace{\kappa\,\|\mathrm{diag}(\rho)-p\|^2}_{\text{consistency}}
+\underbrace{\lambda\,\|\operatorname{skew}(\dot\Theta\Theta^\dagger)\|_F^2}_{\text{gauge}}
\,dt\; }
\]
Here $\|\cdot\|_p$ is the Fisher/Shahshahani norm on $\Delta^\circ$ and $\|\cdot\|_B$ is a Bures‑type norm on density matrices.
Limits: (i) \emph{pure interference}: $\varepsilon^{-1}\!\to 0$, unitary $\mathcal{L}$; (ii) \emph{strong measurement}: large $\gamma(t)$ from witnesses $\Rightarrow$ $Q_F\uparrow$, fast PER commit.

\begin{framed}\small
\textbf{Algorithm (coupled PER/SQC loop).}
At each step $t$: (1) PER: $p\!\leftarrow\!\mathrm{MWU}(p,f)$ with $f=s(t)+K(t)p-\lambda\kappa(t)$; (2) choose $\Gamma$ for that step and $\Theta$ with $|\Theta|^2=\Gamma$ in minimal‑twist gauge; (3) SQC: $\rho\!\leftarrow\!\Theta\rho\Theta^\dagger$ plus Lindblad with rate $\gamma(t)\propto [R_w(t)]_+$; (4) enforce/penalize $\mathrm{diag}\,\rho=p$; (5) compute $R_w,\,R_1,\,Q_F$; (6) stop when the record rule (entropy/concentration + hold) is met.
\end{framed}

% ---------------------------------------------------------
\section{Weak‑measurement view and testable lags}
\label{sec:weak}

\paragraph{Pre/post selection in PER/SQC.}
Let $F$ be the projector for the realized record (Real stage). For any observable $A=\sum_i a_i P_i$ on the SQC layer,
\[
A_w(t) \;=\; \frac{\mathrm{tr}\!\big(F\,A\,\rho(t)\big)}{\mathrm{tr}\!\big(F\,\rho(t)\big)}
\]
is the PER‑adapted weak value (column phases fixed by minimal‑twist gauge). Weak couplings keep $Q_F\!\approx\!0$ and $R_w\!\approx\!0$; strong couplings raise both.

\paragraph{Lag prediction.}
On the \emph{rise} (weak‑information regime) the \emph{witness rate} tends to \emph{lead} the disturbance: $R_w$ leads $\dot Q_F$;
near \emph{collapse} the sign flips (nonunitarity leads).
These are exactly the patterns we observe in LIGO and the Bell data (rise: small positive lags; full‑trace: negative lags), and they co‑localize with spikes in $R_w$, $R_1$, and FDR.

% ---------------------------------------------------------
\section{Interference patterns in PER/SQC (double slit)}
\label{sec:double-slit}

\paragraph{Fringes from SQC; recording from PER.}
Let $\psi(x)=\psi_1(x)+\psi_2(x)$. With which‑path entanglement $|E_1\rangle,|E_2\rangle$ the screen intensity is
\[
I(x;\gamma)=|\psi_1|^2+|\psi_2|^2
+2\,|\gamma|\,\mathrm{Re}\!\big[\psi_1^*(x)\psi_2(x)\big],\qquad
\gamma=\langle E_2|E_1\rangle.
\]
Thus $|\gamma|$ sets visibility (SQC), while PER supplies the stopping rule (click formation).
In Fraunhofer geometry with slit width $a$, separation $d$, distance $L$, wavelength $\lambda$,
\[
I(x;\gamma)\ \propto\ 2\,\mathrm{sinc}^2\!\!\left(\frac{\pi a x}{\lambda L}\right)\!\Big[\,1+|\gamma|\cos\!\Big(\tfrac{2\pi d}{\lambda L}x+\varphi_0\Big)\Big].
\]
(Any AB phase adds to $\varphi_0$.)

\begin{figure}[h]
  \centering
  \includegraphics[width=.68\linewidth]{double_slit_interference_demo.png}
  \caption{\textbf{Double‑slit.} Predicted $I(x;\gamma)$ for $|\gamma|\in\{1,0.6,0\}$.
  PER’s $R_w$ and SQC’s $Q_F$ grow with which‑path strength; visibility falls accordingly.}
  \label{fig:double-slit}
\end{figure}

% ---------------------------------------------------------
\section{Quantum Cheshire Cat in PER/SQC}
\label{sec:qcc}

\paragraph{Setup and weak values.}
On $\mathcal{H}_{\rm path}\!\otimes\!\mathcal{H}_{\rm prop}$ with path $\{|A\rangle,|B\rangle\}$ and property (spin/polarization), take
\[
|\psi\rangle=\frac{|A\rangle+|B\rangle}{\sqrt2}\otimes|{+}\rangle_x,\qquad
|\phi\rangle=\frac{|A\rangle\otimes|{+}\rangle_x+|B\rangle\otimes|{-}\rangle_x}{\sqrt2}.
\]
Then the PER‑adapted weak values
\[
\langle \Pi_A\rangle_w=1,\ \ \langle \Pi_B\rangle_w=0,\qquad
\langle \Pi_A\,\sigma_z\rangle_w=0,\ \ \langle \Pi_B\,\sigma_z\rangle_w=1
\]
exhibit the “cat/grin” separation:
a weak path probe responds in arm $A$ while a weak property probe responds in arm $B$, conditioned on the post‑selection $F$.
In PER/SQC this is a \emph{conditional weak‑response} statement in the evanescent regime; as coupling grows, $Q_F$ rises, the interference dies, and the effect disappears.


% =========================================================

% =========================================================
\section{Causality, coherent broadcasts, and the intrinsic clock}
\label{sec:causality}

\paragraph{Front velocity vs.\ phase/group velocity.}
Only the \emph{front} (signal) velocity $v_f$ matters for causality: it is the speed at which the first nonzero disturbance can arrive. 
Phase or group velocities may exceed $c$ in dispersive media without conveying new information; the causal front obeys $v_f\le c$ (and in vacuum $v_f=c$).

\paragraph{Causal witness kernels.}
We encode witnesses by retarded convolutions
\begin{equation}
s_i(x)\;=\;\int d^4y\;\,G_{\mathrm{ret}}(x-y)\,J_i(y),\qquad x=(\mathbf x,t),
\end{equation}
with $G_{\mathrm{ret}}(x-y)=0$ for spacelike separations $(x-y)^2<0$.
The PER driver $f_i=s_i+(Kp)_i-\lambda\kappa_i$ then yields
$\dot p_i(\mathbf x,t)=0$ until $s_i(\mathbf x,t)\neq 0$.
Hence the \emph{earliest} change of $(p_i,C,\theta)$ at $(\mathbf x,t)$ is light-cone limited by the sources of $J$.

\begin{lemma}[Coherence-front bound]
Let $A=(\mathbf y,t')$ be a source event with compact current $J$, and $B=(\mathbf x,t)$ a detector event.
Assume (i) witnesses use a retarded kernel $G_{\mathrm{ret}}$ supported inside the light cone; (ii) PER updates at $(\mathbf x,t)$ depend only on $s(\mathbf x,t)$.
Then the intrinsic clock $\theta(\mathbf x,t)$ (with $\dot\theta=\alpha\,\dot C$, or its forward-only/smoothed variants) satisfies
\begin{equation}
\dot\theta(\mathbf x,t)=0\quad\text{whenever}\quad t-t' < \frac{\|\mathbf x-\mathbf y\|}{c}\,.
\end{equation}
Equivalently, the set where $\dot\theta$ first becomes positive due to $A$ is contained in $J^+(A)$ (the future light cone of $A$).
\end{lemma}

\noindent\emph{Proof sketch.}
By causality, $s_i(\mathbf x,t)=(G_{\mathrm{ret}}*J_i)(\mathbf x,t)=0$ until $t-t'\ge \|\mathbf x-\mathbf y\|/c$.
Therefore $f_i(\mathbf x,t)$ and $\dot p_i$ vanish; $\dot C$ (a Fisher-quadratic in $f$) remains zero; hence $\dot\theta=\alpha\,\dot C=0$ until the front arrives.

\paragraph{SQC microcausality.}
With a local, gauge-covariant generator $\dot\rho=-i[H(\mathbf A,\phi,g_{\mu\nu}),\rho]+\mathcal D(\rho)$, microcausality ($[\hat{\mathcal O}(x),\hat{\mathcal O}(y)]=0$ for spacelike $x-y$) implies that neither the coherent part nor the GKLS dissipator can alter $\rho$ outside the light cone.
Thus the nonunitarity diagnostic $Q_F$ at $(\mathbf x,t)$ can rise only after the witness front intersects the detector worldline.

\paragraph{Operational speed and coherence cones.}
Define the operational bound for broadcast-driven commits
\begin{equation}
c_{\rm op}\;:=\;\sup\ \frac{\|\mathbf x_B-\mathbf x_A\|}{\,T_{\rm commit}(B)-T_{\rm emit}(A)\,}\,,
\end{equation}
taken over events where the only driver is a broadcast. In vacuum $c_{\rm op}=c$.
In HK (reaction--transport) form, impose $\|\mathbf v(x,t)\|\le c$; PER probability mass then propagates within \emph{coherence cones} analogous to light cones.

\paragraph{Implications for data.}
Our rise-phase lags (e.g.\ LIGO \SIrange{32}{48}{ms} witness-lead) and collapse-phase lags (e.g.\ Bell full-trace collapse-lead) compare local time series \emph{after} front arrival.
They reflect the weak$\to$strong measurement crossover under causal witnesses; they do not imply superluminal propagation.


% =========================================================
\section{Entropy, potentialization, and the force dictionary}
\label{sec:entropy-forces}

\paragraph{Three entropies and ``potentialization.''}
We distinguish: (i) the PER Shannon entropy $H(p)=-\sum_i p_i\ln p_i$ obeying $\dot H=-\mathrm{Cov}_p(f,\ln p)$ and $\dot C=\mathrm{Var}_p(f)\!\ge 0$ for frozen witnesses; (ii) the SQC von Neumann entropy $S(\rho)=-\mathrm{tr}(\rho\ln\rho)$, conserved by unitary transport and increased by Lindbladian decoherence; and (iii) the thermodynamic/environmental entropy $S_{\rm env}$.
We call the rate at which \emph{accessible, redundant} Real records are converted into broadly distributed, non-redundant field degrees of freedom the \emph{potentialization rate} $\Pi(t)$.
Landauer/Clausius considerations give the budget
\begin{framed}\small
\begin{equation}
\frac{d}{dt}\,S_{\rm env}(t)\ \gtrsim\ k_B(\ln 2)\,\Pi(t),
\end{equation}
\end{framed}
with equality in ideal isothermal erasure.
Operationally, potentialization is indicated by \emph{source-side} SQC nonunitarity increases ($Q_F\uparrow$), little immediate decoding by witnesses ($R_w\downarrow$ at receivers), and a drop or stall in local PER coherence accumulation.

\paragraph{Replication vs.\ potentialization (index).}
We define a dimensionless index
\begin{equation}
P \;\equiv\; 1 - \frac{\text{rate of new decoded records (receivers)}}{\text{rate of source record loss}},
\end{equation}
so that $P\!\approx\!1$ indicates near-pure potentialization (e.g.\ thermal radiation far from detectors) and $P\!\approx\!0$ indicates near-pure replication (e.g.\ coherent, structured broadcasts that promptly create new Real at receivers).
In practice the numerator/denominator can be estimated from PER/SQC traces: receiver-side $R_w$ and source-side $Q_F$ (and local $C$) with appropriate coarse-graining.

\paragraph{Forces in PER/SQC (one dictionary).}
We collect conservative, dissipative, and informational drives in the coupled $(\rho,p)$ system:
\begin{center}
\begin{tabular}{@{}p{2.6cm}p{4.7cm}p{4.9cm}p{3.2cm}@{}}
\toprule
\textbf{Force class} & \textbf{SQC generator piece} & \textbf{PER driver piece} & \textbf{Typical signature} \\
\midrule
Conservative (EM, gravity, etc.) & Hamiltonian $H(t)$ with minimal coupling / metric; unitary $\dot\rho=-i[H,\rho]$ & (none directly; enters via what witnesses measure) & $Q_F\!\approx\!0$, $R_w\!\approx\!0$ until readout; interference patterns shaped \\
Measurement / dephasing & Lindblad $\mathcal D=\sum_\alpha \gamma_\alpha(L_\alpha\rho L_\alpha^\dagger-\tfrac12\{L_\alpha^\dagger L_\alpha,\rho\})$ & penalties $\kappa$, reduced effective coupling, or explicit witness scores $s_i$ & $Q_F\uparrow$, visibility$\downarrow$; potentialization if no receivers \\
Which-path / evidence & weak$\to$strong measurement sequence; GKLS with growing rates & Fisher-gradient drive $\dot p_i=p_i(f_i-\bar f)$ with $f_i=s_i+(Kp)_i-\lambda\kappa_i$ & \emph{Rise}: $R_w$ leads $\dot Q_F$; \emph{Collapse}: $\dot Q_F$ leads \\
Radiation damping / friction & amplitude damping channels & increased $\kappa$, reduced $K$ & drains local Real, $Q_F\uparrow$; often potentialization \\
Entropic/free-energy & contact/GENERIC $h=E-TS$ with $\dot h=-\gamma h$ & Fisher free-energy ascent $\mathrm{grad}_p C$ & $R_1$ and FDR spike at bottlenecks \\
\bottomrule
\end{tabular}
\end{center}

\paragraph{Lag/commit signatures.}
The weak-measurement crossover predicts the \emph{rise} ROI to show a small \emph{witness-lead} (positive lag: $R_w$ before $\dot Q_F$), while the \emph{collapse} shows a \emph{nonunitarity-lead} (negative lag).
This matches our LIGO and Bell analyses and provides a practical classification: early replication vs.\ late potentialization.

\paragraph{Relation to the hybrid action.}
All forces above appear in the unified action
\begin{equation}
\mathcal{S}[p,\rho,\Theta]=\!\int\!
\underbrace{\tfrac{1}{2\varepsilon}\big\|\dot p-\mathrm{grad}_p C\big\|_{p}^{2}}_{\text{information / Fisher}}\;
+\;\underbrace{\tfrac{1}{2\varepsilon_q}\big\|\dot\rho-\mathcal{L}(\rho)\big\|_{B}^{2}}_{\text{conservative \& dissipative}}\;
+\;\kappa\,\|\mathrm{diag}(\rho)-p\|^2
+\;\lambda\,\|\mathrm{skew}(\dot\Theta\Theta^\dagger)\|_F^2\,dt,
\end{equation}
with the slow manifold $\mathrm{diag}\,\rho=p$ and $\pi=f-\bar f$ capturing PER/SQC consistency and the replicator momentum, respectively.

\paragraph{Practical recipe.}
To estimate $P$ and potentialization epochs in data: (i) compute source-side $Q_F(t)$ and local $C(t)$; (ii) sum receiver-side $R_w(t)$ over actual detectors; (iii) coarse-grain both with the same window; (iv) define $\mathcal{R}_{\rm lost}=\int (\text{source record loss})\,dt$ and $\mathcal{R}_{\rm created}=\int (\text{receiver decodes})\,dt$ and set $P=1-\mathcal{R}_{\rm created}/\mathcal{R}_{\rm lost}$. High-$P$ intervals indicate potentialization.


\section{Diagnostics: coherence, information, and witnesses}\label{sec:diagnostics}

% \paragraph{Definitions.}
% Let $C_{\rm ext}$ (witness-driven) and $C_{\rm int}$ (internal) sum to $C$.
% Define the witness ratio and fluctuation--dissipation diagnostic:
% \begin{equation}
%   R_w(t)=\frac{\dot C_{\rm ext}(t)}{\dot C_{\rm int}(t)},\qquad
%   \mathrm{FDR}(t)=\frac{\mathrm{Var}_t[\dot C]}{\E_t[\dot C]}.
% \end{equation}
% We compute $R_w$ stably by smoothing $C_{\rm ext}$, differencing, clipping denominators, and low-pass filtering.

\paragraph{Definitions.}
We set
\[
C(p,t)=\underbrace{\sum_i p_i(t)\,s_i(t)}_{C_{\rm ext}(t)}\ +\ \underbrace{p(t)^\top K(t)\,p(t)\ -\ \lambda\sum_i p_i(t)\kappa_i(t)}_{C_{\rm int}(t)}.
\]
Derivatives: apply the \emph{same} low–pass (boxcar or EWMA, window $\Delta$) to $C_{\rm ext}$ and $C_{\rm int}$, central–difference, then compute
\[
R_w=\frac{\dot C_{\rm ext}}{\dot C_{\rm int}+\varepsilon},\qquad \varepsilon\in[10^{-6},10^{-3}] \text{ for scale}.
\]
Defaults: LIGO $\Delta=\SI{100}{ms}$; Bell $\Delta$ equals one witness step (e.g.\ $40$k trials).

\paragraph{Lag statistic and rise/fall split.}
We correlate $R_w$ with $\dot Q_F$ across integer lags; the best lag is the argmax of $|\mathrm{corr}|$.
For transients, we use a derivative-gated ROI, and we split by $\mathrm{sign}(\dot Q_F)$ into \emph{rise} and \emph{fall}.

% =========================================================
\section{Experimental design (common pipeline)}\label{sec:design}

\begin{enumerate}[leftmargin=1.2em, itemsep=2pt]
\item \textbf{Witnesses.}
LIGO: cumulative dissipation and redundancy of local carriers;
Bell: prefix-sum CHSH windows over all four settings, mapping $S(t)$ to two-path supports $(\mathrm{H}_0:\,S\!\le 2,\ \mathrm{H}_1:\,S>2)$ via a cumulative or local (EWMA/boxcar) map.
\item \textbf{PER.} MWU/replicator on supports $\to$ $p(t)$; compute $C,H,R_1,R_w$.
\item \textbf{SQC.} OT/lazy transport on $p(t)$ with entropy $\varepsilon$ $\to$ $Q_F(t)$ (two-path identity cost).
\item \textbf{Alignment \& lags.} Align SQC to the PER grid ($t_{\rm SQC}=t_{1:}$); compute lags on (i) ROI and (ii) full trace; report record times (PER entropy threshold; SQC threshold $\tau$ with hold).
\end{enumerate}

% =========================================================
\section{Results I: LIGO GW150914}\label{sec:gwosc}

We analyze \texttt{GWOSC} strain around GW150914 (H and L sites; standard filtering).
In early evanescence (\emph{rise}), a derivative-gated ROI (z-threshold 0.5, pad 10) yields a short \emph{positive} lag (witness-lead)
\begin{equation*}
\text{lag} \approx \SI{0.032}{s}\ \text{to}\ \SI{0.048}{s}\quad\text{(positive: $R_w$ leads $\dot Q_F$)},
\end{equation*}
with moderate correlations (0.35--0.45).\footnote{Representative ROI metrics: $\{0.032\,\mathrm{s},\,0.28\mbox{--}0.40\}$ and $0.048\,\mathrm{s}$ runs; see the repository figures and JSONs.}
On the full trace (rise+fall) the best lag is \emph{negative} $\approx\SI{-0.112}{s}$ (collapse-lead).
PER commits before SQC (\emph{e.g.} $T_{\rm PER}\approx\SI{15.60}{s}$ vs.\ $T_{\rm SQC}\approx\SI{15.73}{s}$).

\begin{figure}[t]
\centering
\includegraphics[width=.47\linewidth]{gwosc_QF_norm.png}\hfill
\includegraphics[width=.47\linewidth]{gwosc_Rw.png}\\[4pt]
\includegraphics[width=.47\linewidth]{gwosc_lag_scan_roi.png}\hfill
\includegraphics[width=.47\linewidth]{gwosc_overlay.png}
\caption{\textbf{LIGO GW150914.} Top: SQC nonunitarity $\tilde Q_F(t)$ and PER witness ratio $R_w(t)$ (seconds). Bottom: ROI lag scan (positive means $R_w$ leads) and overlay at the best ROI lag. ROI shows a short positive lead ($\sim\!30$--$50$\,ms); full trace flips sign during collapse.}
\label{fig:gwosc}
\end{figure}

\vspace{-8pt}
\paragraph{Numbers used in Fig.~\ref{fig:gwosc}.}
ROI: $+0.032$\,s and $+0.048$\,s best lags (positive) with moderate correlations; full trace: $-0.112$\,s; commits: PER $15.60$\,s vs.\ SQC $15.73$\,s.

% =========================================================
\section{Results II: NIST CH Pockels Bell test (100\,kHz)}\label{sec:bell}

We construct a trial-wise CSV and compute CHSH via prefix sums with windows $(200{,}000,\,40{,}000)$ and $\ge 10$ per cell.
We map $S(t)$ to two-path supports and run PER/SQC as above.

\paragraph{Scale and units.}
After thinning/filters the working set contains $\sim 4.07\times 10^7$ trials; the trial clock is converted at \SI{10}{\micro s}/trial.

\paragraph{Full-trace statistic (conservative supports).}
With a coarse SQC grid ($\mathrm{d}t_{\rm SQC}\!\approx\!\SI{5.078}{s}$) the best lag is
\begin{equation*}
\text{lag}\approx \SI{-30.5}{s}\quad\text{with}\quad \mathrm{corr}\approx 0.316\quad (\dot Q_F\ \text{leads}\ R_w),
\end{equation*}
and with a denser grid ($\mathrm{d}t_{\rm SQC}\!\approx\!\SI{2.536}{s}$) the sign persists with $\text{lag}\approx\SI{-7.61}{s}$.

\paragraph{Early ROI and local witnesses.}
Restricting to the first 10--20\% of the SQC timeline and switching to local supports (EWMA/boxcar) consistently \emph{shrinks} the magnitude and can flip the sign in subsets, but a conservative, full-trace statistic remains \emph{collapse-lead}.

\begin{figure}[t]
\centering
\includegraphics[width=.47\linewidth]{bell_QF_norm.png}\hfill
\includegraphics[width=.47\linewidth]{bell_Rw.png}\\[4pt]
\includegraphics[width=.47\linewidth]{bell_lag_scan_rise.png}\hfill
\includegraphics[width=.47\linewidth]{bell_lag_scan.png}
\caption{\textbf{NIST Bell.} Top: $\tilde Q_F(t)$ and $R_w(t)$ in seconds. Bottom-left: rise-only lag scan (negative indicates $\dot Q_F$ leads); bottom-right: full-trace lag scan. Conservative smoothing emphasizes collapse-lead; local supports and ROI-gating reduce the lag magnitude.}
\label{fig:bell}
\end{figure}

\vspace{-8pt}
\paragraph{Numbers used in Fig.~\ref{fig:bell}.}
Full-trace lags $\{-30.47\,\mathrm{s},\,-7.61\,\mathrm{s}\}$ for $\mathrm{d}t_{\rm SQC}\!\approx\!\{5.078\,\mathrm{s},\,2.536\,\mathrm{s}\}$, respectively; working set size $\sim 4.07\times 10^7$ trials.

% =========================================================
\section{Discussion}
\begin{enumerate}[leftmargin=1.2em]
\item \textbf{Rise vs.\ fall universality.} Both case studies exhibit \emph{collapse-lead} (negative lags) on the fall/whole trace, while the rise often shows a short \emph{witness-lead} once we restrict to the evanescent ROI and use local witnesses.
\item \textbf{PER$\to$SQC handoff.} PER commits slightly before SQC (LIGO), consistent with witness assimilation preceding amplitude-level (near-)unitarity.
\item \textbf{Diagnostics.} $R_w$ and $R_1$ spike at transitions; FDR grows near quasi-critical slowdowns; both support the action picture.
\end{enumerate}

% =========================================================
\section{Methods (key settings)}
Unless otherwise noted:
\begin{itemize}[leftmargin=1.2em]
\item \textbf{PER:} MWU step $\eta\!=\!0.6$, no internal regularizer ($\alpha\!=\!0$), initial $p_0=(\tfrac12,\tfrac12)$.
\item \textbf{Witness supports:} CHSH windows $(200\mathrm{k},40\mathrm{k})$; supports from $S(t)$ via cumulative (conservative) or EWMA (span 4--6) or boxcar (8--10) on an $N_{\rm steps}$ grid (600--800).
\item \textbf{SQC:} OT with identity cost ($N\!=\!2$) and $\varepsilon\!=\!0.06$; SQC commit uses threshold $\tau\!=\!0.20$ (LIGO) and $0.20$--$0.30$ (Bell) with hold $0.20$\,s (LIGO) and $0.20$ grid units (Bell).
\item \textbf{Lag:} z-gated ROI on $|\mathrm{z}(\dot Q_F)|>\!0.5$, pad 10, Pearson correlation over lags $|k|\le L/8$; surrogates: phase or circular shifts; for Bell the trial clock uses \SI{10}{\micro s}/trial.
\end{itemize}

\subsection*{Lag–scan significance and familywise control}
We scan integer lags $k\in[-K,K]$ and use the maximum absolute Pearson correlation as the test statistic.
We build phase–randomized (or circular–shift) surrogates that preserve the spectra and the ROI rule, compute the null distribution of the \emph{max–across–lags} statistic, and report $p$–values and 95\% CIs.
We also report Bartlett–corrected effects after AR(1) prewhitening.

\subsection*{ROI gating and bias control}
We use $|z(\dot Q_F)|>\tau_z$ with padding as default. To test circularity we (i) gate on witness side $|z(R_w)|$, and (ii) use a symmetric union gate.
The chosen gate is preserved in surrogates to estimate any bias.

\subsection*{Alignment and common time base}
We interpolate (or aggregate) both $R_w$ and $Q_F$ to a common, oversampled grid before lag scans; unit conversions (e.g.\ $10\,\mu$s/trial for Bell) are applied once on that grid.

\subsection*{Witness construction for Bell}
Beyond cumulative windows we test EWMA, boxcar, per–setting witnesses, and randomized blocks; across these choices the full–trace sign (collapse–lead) persists in our runs, with magnitude changes reported in the supplement.

\subsection*{Commit thresholds and holds}
We report commit times with the $(\tau,\text{hold})$ used and include a sensitivity sweep in code (typical: $\tau\in[0.1,0.3]$, hold $\in[0.15,0.25]$\,s for LIGO; Bell in grid units).

% =========================================================
\section*{Availability}
Python scaffolds, the LIGO and NIST drivers (including the rise/fall split), and all figures used here come from the same PER/SQC repository companion.

\appendix
\section{Deriving the Shahshahani geometry and geodesics}\label{app:shahshahani}
On $\Delta^\circ=\{p>0:\sum_i p_i=1\}$ with $T_p\Delta=\{u:\sum_i u_i=0\}$, define $G_p(u,v)=\sum_i u_i v_i/p_i$.
The square–root map $F:\Delta^\circ\to\mathbb{S}^{N-1}_2$, $x_i=2\sqrt{p_i}$, pulls back the Euclidean metric to $G$; hence geodesics correspond to great circles under $F$.


\end{document}
